% Ad Atticum 8.8 (ad-atticum-08-08)
% Source: https://raw.githubusercontent.com/PerseusDL/canonical-latinLit/master/data/phi0474/phi057/phi0474.phi057.perseus-lat1.xml
% Fetched by scripts/fetch_latin.py

% Dateline (Perseus): Scr. in Formiano vi K. Mart. a. 705(49).

\ciceroLetterOpener{CICERO ATTICO salutem}

\ciceroSection{1} o rem turpem et ea re miseram! sic enim sentio, id demum aut potius id solum esse miserum quod turpe sit. aluerat Caesarem; eundem repente timere coeperat, condicionem pacis nullam probarat, nihil ad bellum pararat, urbem reliquerat, Picenum amiserat culpa, in Apuliam se compegerat, ibat in Graeciam, omnis nos ἀπροσφωνήτουσ , expertis sui tanti, tam inusitati consili relinquebat.

\ciceroSection{2} ecce subito litterae Domiti ad illum, ipsius ad consules. fulsisse mihi videbatur τὸ καλὸν ad oculos eius et exclamasse ille vir qui esse debuit, πρὸσ ταῦθ' ὅ τι χρὴ καὶ παλαμάσθων καὶ πάντ' ἐπ' ἐμοὶ τεκταινέσθων: τὸ γὰρ εὖ μετ' ἐμοῦ. at ille tibi πολλὰ χαίρειν τῷ καλῷ dicens pergit Brundisium. Domitium autem aiunt re audita et eos qui una essent se tradidisse. o rem lugubrem! itaque intercludor dolore quo minus ad te plura scribam. tuas litteras exspecto.
